\chapter{Gout}

\section*{Definition}
Gout is a crystal arthropathy caused by the deposition of monosodium urate crystals in joints and soft tissues, presenting with acute, self-limited episodes of intense pain, swelling, erythema, and warmth. A serum urate level above 6.8 mg/dL indicates supersaturation, though not all hyperuricemic individuals develop gout.

\section*{What Happens in Your Body}
Chronic hyperuricemia leads to the formation of monosodium urate crystals in synovial fluid. Crystal shedding triggers an innate immune response via NLRP3 inflammasome activation, releasing IL-1$\beta$, IL-6, and TNF-$\alpha$. Neutrophil recruitment and phagocytosis of crystals cause a robust inflammatory reaction. Tophi represent chronic granulomatous deposits of urate crystals surrounded by foreign-body giant cells.

\section*{Causes}
\begin{itemize}
  \item \textbf{Primary (idiopathic):} Uric acid underexcretion or overproduction with no identifiable cause
  \item \textbf{Dietary:} Purine-rich foods (red meat, organ meats, shellfish), fructose-sweetened beverages, alcohol (especially beer and spirits)
  \item \textbf{Medication-related:} Thiazide and loop diuretics, low-dose aspirin, cyclosporine, calcineurin inhibitors
  \item \textbf{Comorbidities:} Chronic kidney disease, obesity, hypertension, diabetes mellitus, metabolic syndrome
  \item \textbf{Genetic:} Urate transporter gene variants (SLC2A9, ABCG2, URAT1)
  \item \textbf{Hematologic:} Tumor lysis syndrome, myeloproliferative disorders, hemolytic anemia, psoriasis
\end{itemize}

\section*{When to See a Doctor}
See your doctor if:
\begin{itemize}
  \item Acute onset of severe, throbbing pain in a single joint (classically the first metatarsophalangeal joint)
  \item Joint erythema, warmth, and swelling with exquisite tenderness
  \item Fever or chills accompanying joint symptoms (rule out septic arthritis)
  \item Recurrent episodes or tophaceous deposits
\end{itemize}

\section*{Related Symptoms}
\begin{itemize}
  \item Podagra: classic first MTP joint involvement
  \item Tophi: subcutaneous urate deposits in ears, elbows, fingers, and Achilles tendon
  \item Joint erythema and desquamation after an acute flare
  \item Renal calculi (urate nephrolithiasis)
  \item Intercritical periods with no symptoms between flares
\end{itemize}
\textbf{Important:} Septic arthritis can coexist with gout; joint aspiration for Gram stain, culture, and crystal analysis is mandatory when infection is suspected.

\section*{Self-Care / Home Management}
\begin{itemize}
  \item Rest and immobilization of the affected joint during acute flares
  \item Ice packs applied for 20 minutes several times daily
  \item Hydration with water to promote urate excretion
  \item Avoidance of alcohol, purine-rich foods, and fructose-sweetened beverages
  \item Over-the-counter NSAIDs (ibuprofen, naproxen) if no contraindications
\end{itemize}

\section*{How Your Doctor Will Evaluate You}
\begin{itemize}
  \item Joint aspiration with polarized light microscopy showing negatively birefringent needle-shaped crystals
  \item Serum urate level (normal does not exclude acute gout; level may be low during a flare)
  \item Synovial fluid WBC count with differential to exclude septic arthritis
  \item Plain radiography: may show punched-out erosions with overhanging edges (rat-bite lesions) in chronic tophaceous gout
  \item Ultrasound for double-contour sign and tophi detection
  \item DECT (dual-energy CT) for urate deposition mapping
\end{itemize}

\section*{Treatment Options}
\begin{itemize}
  \item \textbf{Acute flare:} NSAIDs (indomethacin, naproxen), colchicine (start within 36 hours), or corticosteroids (oral, intra-articular, or IM)
  \item \textbf{Urate-lowering therapy:} Allopurinol or febuxostat as first-line; probenecid for underexcretors
  \item Goal serum urate $<$ 6.0 mg/dL ($<$ 5.0 mg/dL with tophi)
  \item Prophylaxis with low-dose colchicine or NSAIDs during initiation of urate-lowering therapy (3--6 months)
  \item Lifestyle modification: weight loss, dietary counseling, alcohol restriction
  \item Surgical excision of large tophi causing compressive symptoms
\end{itemize}

\section*{References}
\begin{thebibliography}{99}
\bibitem{gout:ref1} Neogi T, Jansen TLTA, Dalbeth N, et al. ``2015 Gout Classification Criteria: An American College of Rheumatology/European League Against Rheumatism Collaborative Initiative.'' Ann Rheum Dis. 2015;74(10):1789--1798.
\bibitem{gout:ref2} Dalbeth N, Merriman TR, Stamp LK. ``Gout.'' Lancet. 2016;388(10055):2039--2052.
\bibitem{gout:ref3} Khanna D, Fitzgerald JD, Khanna PP, et al. ``2012 American College of Rheumatology Guidelines for Management of Gout.'' Arthritis Care Res. 2012;64(10):1431--1446.
\bibitem{gout:ref4} Terkeltaub RA. ``Clinical Practice: Gout.'' N Engl J Med. 2003;349(17):1647--1655.
\bibitem{gout:ref5} Richette P, Bardin T. ``Gout.'' Lancet. 2010;375(9711):318--328.
\bibitem{gout:ref6} Choi HK, Atkinson K, Karlson EW, et al. ``Purine-Rich Foods, Dairy and Protein Intake, and the Risk of Gout in Men.'' N Engl J Med. 2004;350(11):1093--1103.
\end{thebibliography}
